\chapter{Numerical Diagnostic Protocols}
\label{app:numerical-diagnostic-protocols}

This appendix collects the numerical protocols used by
Chapters~\ref{ch:diagnosing-multistability-chaos} and
\ref{ch:research-program}. It is written as an implementation checklist rather
than as a second exposition of the theory. The values in the protocol are
templates for course notebooks unless a specific simulation study supplies
different calibrated values.

\section{Run Manifest}
\label{appC:run-manifest}

Every simulation should produce a manifest with the following fields:
\begin{enumerate}
    \item architecture: E-only, paired E/I, fully connected cluster graph,
    neuron-level Erdos-Renyi graph, or cluster-level Erdos-Renyi graph;
    \item sizes: cluster count \(Q\), local counts \(n_E,n_I\), total counts,
    and the finite-size coordinate used in the plot;
    \item weights: baseline \(J_{\alpha\beta}\), potentiation factors,
    depression factors, fixed indegrees, realized in-degree histograms, and
    heterogeneity amplitudes \(\sigma_{\alpha\beta}^{\mathrm{het}}\);
    \item seeds: connectivity seed, initial-condition seed, and replayable
    stochastic-stream seed;
    \item numerical settings: time step, burn-in, recording length, bin width,
    rate filter, threshold rule, hysteresis width, and minimum dwell time;
    \item outputs: spike file, filtered-rate file, state-transition table,
    current-component file, and perturbation-growth file when available.
\end{enumerate}
The manifest is the unit of reproducibility. A diagnostic plot should be
traceable to a list of manifests, not only to a plotting script.

\section{Network Generation}
\label{appC:network-generation}

The generator should be deterministic given the manifest seed.
\begin{enumerate}
    \item Create labels \((a,\alpha,i)\), where \(a\in\{1,\ldots,Q\}\),
    \(\alpha\in\{E,I\}\), and \(i\) indexes neurons inside a cluster.
    \item For each postsynaptic block \((a,\alpha)\) and presynaptic type
    \(\beta\), construct the within-cluster input set and the outside-cluster
    input sets. For fixed indegree, sample without replacement and verify the
    realized indegree exactly.
    \item Apply within-cluster potentiation factors to the appropriate E/E,
    E/I, I/E, and I/I blocks. Apply compensating depression factors to
    outside-cluster blocks so that the intended mean input is preserved.
    \item Draw quenched outside-cluster factors \(\check j_{ab}^{\alpha\beta}\)
    or adjacency variables. If the model excludes sign-changing weights, record
    whether negative draws are resampled, clipped, or sampled from a truncated
    distribution.
    \item Write block summaries: mean weight, standard deviation, number of
    nonzero entries, and in-degree distribution for each
    \((\alpha,\beta,\mathrm{in/out})\) block.
\end{enumerate}

\section{Spike-to-Rate Pipeline}
\label{appC:spike-rate-pipeline}

For each cluster \(a\) and type \(\alpha\), compute binned activity
\[
    B_{a,m}^{\alpha}
    =
    \frac{1}{n_\alpha\Delta_b}
    \sum_{i=1}^{n_\alpha}
    \int_{m\Delta_b}^{(m+1)\Delta_b}s_{a,i}^{\alpha}(t)\,dt .
\]
Then compute a filtered rate
\[
    r_{a,m}^{\alpha}
    =
    e^{-\Delta_b/\tau_f}r_{a,m-1}^{\alpha}
    +
    \left(1-e^{-\Delta_b/\tau_f}\right)B_{a,m}^{\alpha}.
\]
The first several filter constants after burn-in should be discarded. If an
offline acausal smoother is used for state detection, keep the causal filtered
rate as a separate diagnostic channel so perturbation timing is not blurred by
future data.

\section{Cluster-State Detector}
\label{appC:state-detector}

The detector converts \(r_a^E(t)\) into a binary state vector \(S_a(t)\).
\begin{enumerate}
    \item Pool the post-burn-in rates across clusters and times, then choose
    \(r_0\) from a distribution valley, a fitted mixture model, or a fixed
    quantile specified before the sweep.
    \item Set hysteresis thresholds \(r_{\mathrm{on}}>r_0>r_{\mathrm{off}}\).
    A cluster turns on at an upward crossing of \(r_{\mathrm{on}}\) and turns
    off at a downward crossing of \(r_{\mathrm{off}}\).
    \item Remove events shorter than \(T_{\min}\) and merge interruptions
    shorter than \(T_{\min}\). The value of \(T_{\min}\) should be tied to the
    filter width rather than tuned per trace.
    \item Store transition time, direction, cluster id, pre-state, post-state,
    active-cluster count, and whether the event is isolated.
\end{enumerate}
For sensitivity analysis, rerun the detector over a small grid of thresholds
and dwell times and require the phase label to be stable.

\section{Autocorrelation, HWHM, and Kink}
\label{appC:autocorrelation}

For each filtered rate trace \(r_m\), estimate
\[
    \widehat C[\ell]
    =
    \frac{1}{M-\ell}
    \sum_{m=0}^{M-\ell-1}
    (r_m-\overline r)(r_{m+\ell}-\overline r).
\]
Normalize by \(\widehat C[0]\). The HWHM is the first positive lag at which the
normalized autocorrelation crosses one half, with linear interpolation between
adjacent bins. If no crossing occurs, report a censored lower bound.

For the zero-lag kink, remove the raw same-bin spike-count peak and fit
\[
    \rho(\tau)\approx 1-k_\xi|\tau|+\frac{1}{2}b\tau^2
\]
on several small windows \(0<|\tau|\le w\). The finite-size contribution
should shrink with \(n_E\). A broad autocorrelation that remains smooth at zero
lag and persists as \(n_E\) grows is more consistent with deterministic
mesoscopic dynamics than with finite-size hopping.

\section{Transition-Triggered Analysis}
\label{appC:transition-triggered}

For every isolated event \(t_m\), align traces on \(t_m\) and compute:
\begin{enumerate}
    \item mean cluster-rate trajectory and active-count trajectory;
    \item mean current components
    \(I_{\mathrm{in}}^E,I_{\mathrm{in}}^I,I_{\mathrm{out}}^E,I_{\mathrm{out}}^I\);
    \item spatial variances of the same four current components across
    excitatory neurons in the focus cluster;
    \item Fano factors in sliding windows around the transition, with
    rate-matched non-transition control windows;
    \item raw and residual participation-ratio dimensionality of the
    transition cloud.
\end{enumerate}
Conditioning matters. Off-to-on and on-to-off transitions should be separated,
and a specific state transition \(A\to B\) should be analyzed separately when
there are enough events. The directional coordinate is
\[
    z_{\parallel}(t)
    =
    \frac{(r_B-r_A)^{\mathsf T}(r(t)-r_A)}{\|r_B-r_A\|}.
\]

\section{Lyapunov Protocols}
\label{appC:lyapunov}

For a deterministic mesoscopic model \(\dot x=F(x)\), use the tangent equation
\(\dot U=D F(x)U\). At times separated by \(\Delta_R\), factor \(U=QR\),
replace \(U\leftarrow Q\), and accumulate
\[
    \widehat\lambda_j
    =
    \frac{1}{M\Delta_R}\sum_{m=1}^M\log |R_{jj}^{(m)}|.
\]
For only the largest exponent, a single perturbation vector and repeated
renormalization are sufficient. Convergence checks should vary time step,
\(\Delta_R\), norm, transient length, and perturbation dimension.

For spiking simulations, use common-noise perturbations:
\begin{enumerate}
    \item replay the same connectivity and stochastic streams in two copies;
    \item clone the full state after burn-in;
    \item perturb one membrane potential, all neurons in one cluster, one
    spike, or all neurons in one cluster by one spike;
    \item measure filtered-rate distance before saturation;
    \item repeat across perturbation amplitudes and local population sizes.
\end{enumerate}
Independent-noise divergence is a control for stochastic variability, not a
Lyapunov exponent.

\section{Decision Criteria}
\label{appC:decision-criteria}

A parameter point is classified as finite-size hopping when switch times grow
rapidly with \(n_E\), perturbations recover inside a state, the kink amplitude
shrinks with size, and deterministic or common-noise Lyapunov estimates are not
positive in the large-cluster limit. A parameter point is classified as
mesoscopic-chaos evidence when switch times plateau as \(n_E\) grows, small
perturbations separate, the broad autocorrelation persists without a
finite-size kink, the Lyapunov estimate is positive with convergence checks,
and the behavior appears in a deterministic mesoscopic reduction or survives
the \(\xi\to 0\) scaling. Mixed cases should be reported as mixed; they are
often the scientifically useful boundary of the phase diagram.
